\section{Estimação de pose do alvo e percepção relativa em operações de proximidade}
\label{sec_estimacao_pose_alvo_percepcao_relativa}

Nesta seção é apresentada uma revisão da literatura relacionada à estimação de pose do alvo e à percepção relativa em operações de proximidade orbital. São destacadas as referências que tratam do uso de sensores ópticos e ativos, da reconstrução do estado relativo do alvo e da integração entre navegação relativa, guiamento e controle em missões de \gls{rvdb} com alvos cooperativos e não cooperativos.

Com a evolução das missões de serviço em órbita, remoção ativa de detritos e inspeção de satélites, a percepção relativa passou a ocupar papel central nas operações de proximidade. Nessas missões, tornou-se necessário estimar não apenas a posição relativa entre perseguidor e alvo, mas também sua atitude, velocidade relativa, velocidade angular e, em alguns casos, a própria geometria do alvo. Dessa forma, a literatura passou a explorar diferentes combinações de sensores e algoritmos para a reconstrução do estado relativo durante as fases de aproximação curta.

\cite{volpe2018passive} apresentam um método para determinação da pose e da forma de um satélite não cooperativo e possivelmente desconhecido a partir de imagens obtidas por uma câmera passiva, complementadas por um sensor de distância embarcado na espaçonave perseguidora. Os autores formulam um filtro capaz de estimar posição relativa, atitude, velocidade, velocidade angular e geometria aproximada do alvo com base no rastreamento de características em imagens sucessivas. O trabalho mostra que a navegação relativa pode ser realizada sem depender de marcas cooperativas previamente instaladas no alvo, explorando diretamente informações extraídas da cena observada. Os resultados apresentados indicam desempenho satisfatório na estimação do estado relativo e da forma aproximada do alvo.

Em continuidade a essa linha, \cite{volpe2022gnc} apresentam uma arquitetura integrada de guiamento, navegação e controle para \gls{rvdb} com um alvo não cooperativo em rotação, utilizando câmera monocular passiva. Nesse trabalho, a estimação de pose não é tratada isoladamente, mas integrada a uma arquitetura fechada de navegação e controle. Os autores utilizam extração e rastreamento de características visuais, em conjunto com um filtro \gls{ukf}, para estimar posição relativa, velocidade, atitude, velocidade angular e forma reconstruída do alvo. A partir dessas estimativas, é gerada uma trajetória ótima livre de colisão, validada em co-simulação com imagens renderizadas em malha fechada. Essa referência mostra uma evolução da literatura no sentido da integração entre percepção relativa e decisão de aproximação.

Além das abordagens baseadas em sensores passivos, a literatura recente também apresenta o uso de sensores ativos para contornar limitações associadas à iluminação, à baixa textura superficial do alvo e às distorções provocadas pelo movimento relativo. Nesse contexto, \cite{sun2022tof} apresentam um sistema de estimação de pose de alvos espaciais não cooperativos baseado em câmera \gls{tof}.
Os autores desenvolvem uma arquitetura composta por pré-processamento dos dados, registro quadro a quadro e atualização de modelo global, com o objetivo de reduzir o erro acumulado e o custo computacional. O método é validado por experimentos semi-físicos, apresentando erro translacional mínimo de 8~mm, erro angular de aproximadamente 1$^\circ$, tempo de processamento em torno de 100~ms e frequência de rastreamento de 10~Hz. Esses resultados mostram a viabilidade do uso de sensores ativos compactos em operações de proximidade.

Mais recentemente, \cite{gavilanez2024vision} apresentam uma abordagem para estimação relativa de posição e atitude de espaçonaves não cooperativas baseada em visão estéreo e arquitetura generativa. O método combina extração de pontos característicos, reconstrução de nuvem de pontos por rede generativa adversarial e alinhamento por meio do algoritmo \gls{gicp}.
Os autores buscam melhorar a resolução e a distribuição das nuvens de pontos antes da etapa de registro, com o objetivo de aumentar a precisão final da estimação. Os resultados experimentais em bancada indicam melhora no desempenho da determinação de posição e atitude, sugerindo que técnicas de aprendizado de máquina podem ampliar a robustez de métodos visuais em operações de proximidade.

Paralelamente ao desenvolvimento de métodos específicos, surgiram trabalhos de revisão dedicados à sistematização da navegação relativa baseada em aprendizado profundo. \cite{song2022survey} apresentam uma revisão de métodos de navegação relativa de espaçonaves baseados em \textit{deep learning}, discutindo aplicações em \gls{rvdb}, exploração de pequenos corpos e navegação por terreno, além de bases de dados e desafios de implementação embarcada. Por sua vez, \cite{pauly2023survey} apresentam uma revisão voltada especificamente à estimação monocular de pose de espaçonaves, destacando limitações importantes da literatura recente, como a dificuldade de implantação de modelos computacionalmente intensivos em hardware espacial e a redução de desempenho quando métodos treinados com imagens sintéticas são aplicados a dados mais próximos de condições reais. Essas revisões mostram que, embora o uso de aprendizado profundo tenha ampliado significativamente o repertório técnico da área, ainda permanecem desafios relacionados à confiabilidade, à representatividade dos dados e à transferência entre ambiente simulado e operação real.

De modo geral, a literatura revisada mostra que a estimação de pose do alvo e a percepção relativa passaram a constituir elementos centrais das operações autônomas de proximidade. Enquanto trabalhos anteriores enfatizavam a reconstrução geométrica e a obtenção do estado relativo por sensores ópticos, estudos mais recentes passaram a incorporar integração com arquiteturas de guiamento, navegação e controle, uso de sensores ativos, técnicas baseadas em aprendizado de máquina e preocupações explícitas com robustez computacional e implementação embarcada. Esses aspectos são relevantes para a presente dissertação, pois mostram que a fase de aproximação curta depende não apenas da modelagem da dinâmica relativa, mas também da disponibilidade de informações confiáveis sobre a configuração espacial do alvo ao longo da operação.